\chapter{Fazit und Ausblick}
\label{ch:fazit}

Dieses Kapitel fasst die Ergebnisse entlang der drei Teilfragen zusammen
und leitet daraus eine Empfehlung für die Praxis ab. Der anschließende
Ausblick benennt die Ansatzpunkte, an denen weitere Arbeiten anknüpfen
können.

\section{Fazit}
\label{sec:fazit}

Gegenstand dieser Arbeit war die Frage, wie sich die Netzwerktopologie
auf RL-trainierte Verteidigungsagenten auswirkt (\Cref{sec:ziel}). 
Grundlage waren 120 Trainingsläufe über 24 Matchups aus
vier verteidigten Topologien und sechs Angreifern sowie eine Evaluation mit
9000 Episoden gegen drei Vergleichsstufen (\Cref{ch:ergebnisse}). Für die drei Teilfragen ergibt sich ein unterschiedliches Bild.

Ein Einfluss der Topologie auf die \emph{Lerngeschwindigkeit} ließ
sich nicht nachweisen. Alle 120 Läufe erfüllen das Konvergenzkriterium, im Median
in Episode~42, und die Mediane der vier Topologien liegen mit 41 bis 42,5 so
dicht beieinander, dass sich darauf keine Aussage stützen lässt
(\Cref{sec:konvergenzergebnis}). Das gilt auch für das nachträglich
entworfene Alternativmaß; bei fünf Seeds je Zelle bleibt ein etwaiger
Unterschied unterhalb der Nachweisgrenze (\Cref{sec:limitationen}). Sichtbar
wurde stattdessen eine Eigenschaft der Matchups: Gegen den
Angreifer, der auf der verteidigten Topologie trainiert wurde, dauert
es fast durchweg am längsten, bis ein Lauf konvergiert.

Deutlich trennt die Topologie dagegen das \emph{erreichte Ergebnis}. Das
Training wirkt zunächst überall: Der Restanteil des trainierten Modells
liegt in jeder der 24 Zellen unter dem des Zufallsverteidigers
(\Cref{sec:wirksamkeitergebnis}). Wie viel dem Angreifer dennoch gelingt,
bestimmt maßgeblich die Struktur: Mit zunehmender Segmentierung sinkt die
Crown-Jewel-Quote, vom nahezu ungebremsten Zugriff im vollvermaschten
\emph{flat} über die beiden teilsegmentierten Topologien, die sich
untereinander kaum trennen, bis zu \emph{micro\_segmented}, wo der Zugriff
fast vollständig unterbunden ist. Die Auswertung des Aktionsverhaltens
liefert dazu den Mechanismus (\Cref{sec:aktionen}): Wo die Struktur
Engpässe bietet, werden Sperren wirksam; im vollvermaschten Netz weicht
der Angreifer jeder Sperre aus, und auch dem trainierten Verteidiger
bleibt kein greifendes Mittel. Die wirksamsten Verteidiger dieser
Untersuchung sperren kaum, sondern setzen kompromittierte
Knoten konsequent neu auf (\Cref{sec:rewarddiskussion}). Dieser Befund beruht
allerdings auf einer Vereinfachung der Simulation: Der Verteidiger kennt
dort den Zustand jedes Knotens exakt und kann gezielt neu aufsetzen. In einem realen Netz
ist gerade das ungewiss; dort wäre vermutlich die Sperre die naheliegende erste Maßnahme und
das Wiederherstellen aus einem Backup die zweite.

Neben der verteidigten Topologie prägt auch die \emph{Herkunft des
Angreifers} das Ergebnis (\Cref{sec:angreiferherkunft}). Der auf der verteidigten Topologie
trainierte Angreifer beginnt um den Faktor 1,3 bis 4 stärker als jeder
andere und beschreibt damit den ungünstigsten Fall, den ein Verteidiger zu
erwarten hat. Der über alle Topologien gestuft
trainierte Generalist ist den Spezialisten fast überall unterlegen; nur in
\emph{micro\_segmented} kommt er als einziger noch durch, wobei offen
bleibt, ob dafür die Breite seines Trainings oder dessen letzte Station
verantwortlich ist.

Neben diesen inhaltlichen Antworten steht ein methodisches Ergebnis. Der
Reward allein erwies sich als unzureichendes Maß für
Verteidigungsleistung: In \emph{micro\_segmented} fällt die
Crown-Jewel-Quote zwischen Episode~150 und 230 von 38 auf 5~Prozent,
während sich der Reward kaum noch bewegt (\Cref{sec:cj-verlauf}) -- erst
die Zielgrößen Crown-Jewel-Quote und gehaltene Knoten machen sichtbar, was
die Verteidigung tatsächlich leistet. Auch das eigens entworfene
Konvergenzkriterium trennte die Topologien nicht, weil seine Normierung
von den extremen Startwerten dominiert wird; beides ist für nachfolgende
Arbeiten mit \ac{CBS} unmittelbar verwertbar.

Für die Praxis lässt sich daraus -- im Rahmen des Modells und mit den
Einschränkungen aus \Cref{sec:limitationen} -- eine Richtung ableiten:
möglichst konsequente Segmentierung. Sie wirkt in dieser Untersuchung
doppelt. Zum einen begrenzt die Struktur selbst, was ein Angreifer
überhaupt erreichen kann, noch bevor ein Verteidiger handelt: In
\emph{micro\_segmented} bleiben ihm nur 145 der 235 Punkte, und den
DMZ-Angreifer hielt dort bereits der Zufallsverteidiger in jeder Episode
vom Crown Jewel fern (\Cref{tab:cj-eval}). Zum anderen ist Segmentierung die Voraussetzung
dafür, dass aktive Verteidigung überhaupt greift: Erst Engpässe machen
Sperren treffsicher, während im flachen Netz selbst der trainierte Agent
das wertvollste System praktisch nie schützen konnte. Einem Unternehmen,
das automatisierte Abwehr einsetzen will, wäre demnach zu raten, die
Topologie nicht als gegebene Randbedingung zu behandeln, sondern als
ersten Teil der Verteidigung -- die Struktur dämmt den Angreifer ein,
bevor der erste Agent eine Aktion wählt.

\section{Ausblick}
\label{sec:ausblick}

Die größte Einzeleinschränkung dieser Arbeit ist der eingefrorene
Angreifer: Die gemessene Wirksamkeit gilt gegen ein festes Verhalten, auf
das sich der Verteidiger ungestört einspielen konnte
(\Cref{sec:limitationen}). Zwei Erweiterungen lägen nahe: zunächst
Angreifer, die gegen einen zufällig handelnden Verteidiger trainiert
wurden und Gegenwehr damit überhaupt kennen, und darauf aufbauend ein
Versuch, in dem beide Agenten von null an gleichzeitig lernen. Erst dort
zeigt sich, ob die Neutralisierung des passgenauen Angreifers Bestand hat
oder ob der Einfluss der Angreiferherkunft gegen einen ausweichenden
Gegner wächst.

Am gestuften Training des Generalisten blieb unentschieden, ob seine
Stärke auf \emph{micro\_segmented} Breite oder Aktualität ist. Läufe mit
permutierter Stationenfolge würden beides trennen; ergänzend ließe sich
eine Effektivitätsmatrix eingefrorener Angreifer gegen eingefrorene
Verteidiger aller Stationen aufnehmen. Ebenso stehen die in
\Cref{sec:limitationen} benannten Entwurfsprüfungen aus: eine systematisch
variierte Abwehrbonusrate, der Vergleich port- gegen knotenweiser Sperren,
ein Lauf ohne Sperren in \emph{flat} sowie eine ausdrückliche
Nichtstun-Aktion beziehungsweise ein Action-Masking. Ob die ungültigen
Aktionen, zuletzt 62~Prozent aller Verteidigeraktionen, tatsächlich
gewollte Untätigkeit sind oder der Agent schlicht nicht gelernt hat,
welche Aktionen gültig sind, ließe sich erst damit unterscheiden.

Die statistische Basis ließe sich mit mehr Seeds je Zelle und
Konfidenzbereichen verbreitern, wie von Henderson et~al.\
empfohlen~\cite{henderson2019deepreinforcementlearningmatters}; kleinere
Unterschiede zwischen den segmentierten Topologien wären erst damit
aussagekräftig. Für die Lerngeschwindigkeit fehlt zudem weiterhin ein
Maß, das späte Verbesserungen von früher Ruhe trennt; die in
\Cref{sec:limitationen} vorgeschlagene 90-Prozent-Marke wäre ein
Ausgangspunkt.

Schließlich bleibt die Übertragung. \ac{CBS} abstrahiert von
Zeitverhalten, Erkennungsunsicherheit und unvollständiger Sicht des
Verteidigers; die hier gefundene Rangfolge der Topologien ist eine Aussage
über dieses Modell. Größere und heterogenere Netze, eine Sterntopologie
mit gehärtetem, nicht übernehmbarem Zentrum sowie Umgebungen mit
Teilbeobachtbarkeit würden zeigen, ob der doppelte Nutzen der
Segmentierung auch dort trägt, wo reale Netze weniger ideal geschnitten
sind.
